\section{Introduction}
\label{sec:intro}

When a language model reads the tokens \texttt{[``Compu'', ``Ser'', ``ve'']}, it does not process three independent subwords---it recognizes a single entity.
This observation points to a fundamental gap between how LLMs \emph{tokenize} text and how they \emph{represent} it internally.
While byte-pair encoding (BPE) tokenizers split text into a fixed vocabulary of subword units, the model's hidden states appear to merge certain multi-token sequences into unified representations, effectively building an \newterm{implicit vocabulary} that extends far beyond the explicit tokenizer~\cite{feucht2024token}.

Understanding this implicit vocabulary matters for at least three reasons.
First, it provides a window into what semantic units the model actually operates over, which is central to \emph{interpretability}~\cite{geva2020transformer}.
Second, it can inform \emph{tokenizer design}: if models consistently merge certain multi-token sequences, a better tokenizer could capture them explicitly, reducing the representation gap~\cite{yang2024rethinking}.
Third, it connects to the long-standing question of how neural models handle \emph{non-compositional language}---idioms, named entities, and technical terms whose meaning cannot be derived from their parts~\cite{garcia2021probing}.

\citet{feucht2024token} took an important first step by identifying implicit vocabulary items through ``token erasure,'' showing that LLMs rapidly forget token-level information at the last position of multi-token words and entities.
However, their method has two key limitations: it requires pre-trained linear probes specific to each model (available only for Llama-2-7b and Llama-3-8B), and it does not characterize the \emph{compositionality} of the recovered sequences.
Not all multi-token sequences that the model merges are non-compositional---``the cat'' may be merged for efficiency without being a semantic unit---and distinguishing these cases is essential for understanding what the implicit vocabulary actually contains.

We address both limitations.
We introduce a \emph{probe-free} method that measures representation shift, compositional deviation, and internal similarity directly from the model's hidden states, requiring no pre-trained probes and generalizing to any transformer.
We apply this method to \mistral on 150 Wikipedia articles, extracting 6{,}494 unique multi-token implicit vocabulary items.
We then classify a stratified sample of 500 items using GPT-4.1 to obtain compositionality ratings and semantic categories.

Our main finding is that items with higher erasure-like scores are significantly less compositional (Spearman $\rho = -0.182$, $p = 0.003$; Cohen's $d = -0.60$ when comparing high- vs.\ low-scoring items).
Named entities dominate the implicit vocabulary (38\% of meaningful items), followed by technical terms (15\%), fixed expressions (12\%), and compound words (6\%).
The compositionality distribution is continuous, not binary, with technical terms and fixed expressions occupying a middle ground between fully opaque entities and fully transparent phrases.

Our contributions are:
\begin{itemize}[leftmargin=*,itemsep=0pt,topsep=0pt]
    \item We propose a probe-free method for extracting implicit vocabulary items from any transformer, combining representation shift, compositional deviation, and internal similarity signals.
    \item We provide the first systematic characterization of the compositionality spectrum in LLM implicit vocabularies, showing that higher erasure-like scores correlate with lower compositionality.
    \item We analyze the semantic composition of the implicit vocabulary, finding that named entities, technical terms, and fixed expressions---predominantly non-compositional categories---account for the majority of meaningful items.
\end{itemize}

The rest of this paper is organized as follows.
\Secref{sec:related} reviews related work on token erasure, compositionality, and multiword expression processing.
\Secref{sec:method} describes our probe-free extraction method and compositionality classification pipeline.
\Secref{sec:results} presents our experimental results.
\Secref{sec:discussion} discusses implications and limitations, and \secref{sec:conclusion} concludes.
